\chapter{Conclusion}
\label{ch:conclusion}

This thesis explored the field of volumetric path tracing and methods for improving
the performance of volumetric rendering. We examined how the field developed and
what the latest advancements were. The foundation for our work was the Photon Field
Networks paper by Bauer \emph{et~al.}~\cite{bauer_photon_2024}. We extended their
approach by applying it to a path tracer. To do this, we built on an existing
browser-based volumetric path tracer, VPT. We split the radiance calculation into
direct and indirect illumination and made it possible to use the indirect radiance
values as ground truth for our model training. To enable online training, we sent
data from the browser to a training server running on the same machine. During
training, path tracing was already producing images, and after training was
complete, neural rendering allowed us to replace the time-consuming indirect
illumination stage with a faster model inference.

We showed that our method achieved an overall speedup of $2$ in the indirect
illumination \ac{gpu} stage, with the overall frame time speedup dropping to
around $1.2$. Because our model only approximates transparency, we do not reach
path tracing quality in image metrics for converged images, but our method
performs better on initial frames. Our images tend to achieve better metrics
after $0.5$ seconds, and we also showed that a single frame with multiple
samples is less noisy in our case. Because we used path tracing, we were able
to apply a filtering method on top of our ground truth data, which our model
was able to learn. However, we also found limitations of our method in
configurations where it fails to generate a good image, or any image at all.

\section{Future work}

In our work, we focused on demonstrating that our approach functions correctly,
but there are still opportunities to improve and extend it. The model’s forward
pass in the WebGPU shader could be further optimized using different techniques,
and support for WebGPU itself is continuously improving. With better support for
\ac{gpu} computation in the browser, training could also be performed there,
eliminating the need for a separate training server.

As shown in our results, image quality is the main area that could be improved.
We only predicted color and approximated transparency, which led to lower quality
metrics. The model could be extended to predict transparency directly, and we
could experiment with colors in different color spaces. Using colors before gamma
correction might be beneficial for predicting extremely dark scenes, which we
are currently unable to capture. Another limitation is that changing rendering
parameters requires retraining the model from scratch because it is tied to a
specific configuration. This could be addressed by fine-tuning the model instead
of retraining it for small parameter changes, making subsequent training faster.
